% =====================================================================
%  chunks/intro.tex —— 原书 p.4–5（PDF 11–12）
%  Introduction + Introduction to the English Translation
% =====================================================================
% [原文 p.4]
\chapter*{导言}
\addcontentsline{toc}{chapter}{导言}

在本篇中，我们尝试给出拓扑学的轮廓，并对各种拓扑概念略作说明。必须首先说明：除了必要的最低限度之外，本综述的研究对象并不包含称为「一般拓扑学」的那门分支学科——即在集合论与一般范畴论的背景下考察的一般空间及其映射的理论。（毫无疑问，这一主题将由他人详细综述。）有了这一限定，才可以说：本综述所处理的「拓扑学」，正是 19 世纪末被称为 \emph{Analysis Situs}（位置分析）的那门数学学科。它在后来的不同时期逐渐分化出若干分支学科：「组合拓扑学」、「代数拓扑学」、「微分（或光滑）拓扑学」、「同伦论」、「几何拓扑学」。

% [Tutor Note: 作者在这里划定的边界只是「本综述不系统讲一般拓扑学」，并非说全书不用一般拓扑。第 2 章 §1 恰恰从开集、紧性、Hausdorff 性、度量空间这些一般拓扑的观察与术语讲起，作为后面同伦论的语言准备；读第 2 章时不要把这一节当成「跑题」。]

随着拓扑学在数学其他领域中的应用长期增长，以下更多分支学科逐渐定型：整体变分法、整体几何、李群与齐性空间的拓扑学、复流形与代数簇的拓扑学。此外还有动力系统与叶状结构的定性（拓扑）理论、椭圆型与双曲型偏微分方程的拓扑学。最后，在 70 与 80 年代，拓扑方法被成批地应用于现代物理的问题；事实上在若干情形下，若没有拓扑学概念的帮助，就根本无法理解所论真实物理现象的本质。

既然不可能在本综述中逐一处理上述全部主题，我们在这里只能满足于如下一般性的评论。拓扑学在非常广泛的问题上找到了令人印象深刻的应用——这些问题涉及数学对象与物理对象的定性性质与稳定性。与拓扑学一同发展起来的代数工具，则导致了整个现代代数学的重新定向。

近年来的进展表明，现代李群理论及其表示论连同代数几何，都是在源于拓扑学的一组深刻代数观念的基础上达到今日发展水平的。不过，这些学科在应用中所起的作用相当不同：它们主要被用于对具有深刻内在代数对称性的系统作精确的公式化研究。事实上，这一点在早先关于经典力学与数学物理问题的精确解中就已显现；但它只是在现代对「在某种确切意义下可积」的系统的研究中才变得毫无疑义。只需回忆一些例子：逆散射方法、非线性场系统的（代数几何的）有限隙积分、统计物理与量子场论模型的著名解、自对偶规范场，以及弦论。（不过，这一局面的一个方面值得注意：现代代数几何需要一次认真的「有效化」，使其在精神上回到 19 世纪——那时代数几何被视为公式化分析的一部分。）

% [原文 p.5]
本综述是后续一系列拓扑学论著的总导引；在这些论著中，拓扑学的各分支学科的发展将得到更详细的阐述。

\chapter*{英译本导言}
\addcontentsline{toc}{chapter}{英译本导言}

本综述写于 1983–84 年，1986 年（以俄文）出版。英译工作始于 1993 年。鉴于过去十年间拓扑学中出现了若干重要的新思想，我增补了一个附录来概述其中的一些思想，并添加了若干脚注，以使本综述更贴近当前状态。

我对若干人士为本书作出的宝贵贡献表示感谢。M. Stanko 为俄文版承担了繁重的编辑工作；B. Botvinnik 作为英文版的学术编辑做了细致的工作，尤其是在使文本现代化方面；R. Burns 以很高的速度完成了很好的英译；C. Shochet 在马里兰大学为翻译与文本现代化提供建议和帮助。我也感谢其他同事在文本现代化方面的帮助。

\vspace{1.2em}
\begin{flushright}
Sergei P. Novikov\\
1995 年 11 月
\end{flushright}

% CHUNK intro DONE
